\subsection{Exact spectra and trace spread of the original four relation windows}\label{endpoint-window_review--exact-spectra-and-trace-spread-of-the-original-four-relation-windows}

13 September 2026. Independent mathematical continuation of the original arithmetic determinant transport. This note leaves the sealed transport TeX and its PDF unchanged. It derives the complete spectra of the actual window operators, proves the spectral trace extremum directly, and integrates it against the original relative Gram eigenvalues. The packet, tensor degree, polynomial coordinate, measures, and all four cutoff indices remain fixed.

\subsubsection{1. Original common source and projections}\label{endpoint-window_review--original-common-source-and-projections}

Fix the original actual packet ENDPOINTSOURCEB2396A45C7A5608906169E42SLOT0END, tensor degree ENDPOINTSOURCEB2396A45C7A5608906169E42SLOT1END, and its literal cyclic annihilator ENDPOINTSOURCEB2396A45C7A5608906169E42SLOT2END, of degree ENDPOINTSOURCEB2396A45C7A5608906169E42SLOT3END. Keep the common coefficient space

ENDPOINTSOURCEB2396A45C7A5608906169E42SLOT4END

with its original increasing-power frame. Let ENDPOINTSOURCEB2396A45C7A5608906169E42SLOT5END and ENDPOINTSOURCEB2396A45C7A5608906169E42SLOT6END be the two actual positive source Grams at this cutoff, retaining the Gamma mass and arithmetic mass. Set

ENDPOINTSOURCEB2396A45C7A5608906169E42SLOT7END

No rescaling of either Gram is performed. Positivity of ENDPOINTSOURCEB2396A45C7A5608906169E42SLOT8END follows because its quadratic form on a nonzero vector is a convex combination of two strictly positive numbers.

Write ENDPOINTSOURCEB2396A45C7A5608906169E42SLOT9END for the ENDPOINTSOURCEB2396A45C7A5608906169E42SLOT10END-orthogonal projection onto the original coefficient subspace ENDPOINTSOURCEB2396A45C7A5608906169E42SLOT11END. Write ENDPOINTSOURCEB2396A45C7A5608906169E42SLOT12END for the same-metric orthogonal projection onto

ENDPOINTSOURCEB2396A45C7A5608906169E42SLOT13END

Multiplication by the nonzero polynomial ENDPOINTSOURCEB2396A45C7A5608906169E42SLOT14END is injective, so ENDPOINTSOURCEB2396A45C7A5608906169E42SLOT15END. In particular ENDPOINTSOURCEB2396A45C7A5608906169E42SLOT16END. These are the original relation images, with every power of ENDPOINTSOURCEB2396A45C7A5608906169E42SLOT17END and every finite coefficient retained.

The two actual operators are

ENDPOINTSOURCEB2396A45C7A5608906169E42SLOT18END

All projections in WS4 are computed in the one space WS1 with the same form ENDPOINTSOURCEB2396A45C7A5608906169E42SLOT19END. Every equation below is pointwise in this retained parameter unless an integral is displayed.

For completeness, if ENDPOINTSOURCEB2396A45C7A5608906169E42SLOT20END are subspaces and ENDPOINTSOURCEB2396A45C7A5608906169E42SLOT21END their orthogonal projections for that form, then ENDPOINTSOURCEB2396A45C7A5608906169E42SLOT22END because ENDPOINTSOURCEB2396A45C7A5608906169E42SLOT23END. Also ENDPOINTSOURCEB2396A45C7A5608906169E42SLOT24END, because ENDPOINTSOURCEB2396A45C7A5608906169E42SLOT25END. Hence ENDPOINTSOURCEB2396A45C7A5608906169E42SLOT26END is self-adjoint, idempotent, and has image ENDPOINTSOURCEB2396A45C7A5608906169E42SLOT27END. This proves every nested-projection decomposition used below; none of them assumes unrelated subspaces are orthogonal.

\subsubsection{2. The exact relation and polynomial window spectra}\label{endpoint-window_review--the-exact-relation-and-polynomial-window-spectra}

The original relation spaces have the nested flag

ENDPOINTSOURCEB2396A45C7A5608906169E42SLOT28END

of dimensions ENDPOINTSOURCEB2396A45C7A5608906169E42SLOT29END. For ENDPOINTSOURCEB2396A45C7A5608906169E42SLOT30END, the first two spaces coincide; the following middle complement is then the zero space. Define

ENDPOINTSOURCEB2396A45C7A5608906169E42SLOT31END

The four spaces are mutually orthogonal and their dimensions add to ENDPOINTSOURCEB2396A45C7A5608906169E42SLOT32END. The flag decomposition proves that they span ENDPOINTSOURCEB2396A45C7A5608906169E42SLOT33END. On them the three summands in ENDPOINTSOURCEB2396A45C7A5608906169E42SLOT34END act respectively as follows:

ENDPOINTSOURCEB2396A45C7A5608906169E42SLOT35END

Every entry is a restriction of the actual projection, not an assigned new operator.

The original polynomial flag is

ENDPOINTSOURCEB2396A45C7A5608906169E42SLOT36END

Its dimensions are ENDPOINTSOURCEB2396A45C7A5608906169E42SLOT37END; the middle two coincide for ENDPOINTSOURCEB2396A45C7A5608906169E42SLOT38END. The corresponding four orthogonal pieces are

ENDPOINTSOURCEB2396A45C7A5608906169E42SLOT39END

On these pieces ENDPOINTSOURCEB2396A45C7A5608906169E42SLOT40END has respective values ENDPOINTSOURCEB2396A45C7A5608906169E42SLOT41END, while ENDPOINTSOURCEB2396A45C7A5608906169E42SLOT42END has respective values ENDPOINTSOURCEB2396A45C7A5608906169E42SLOT43END. Therefore ENDPOINTSOURCEB2396A45C7A5608906169E42SLOT44END acts by ENDPOINTSOURCEB2396A45C7A5608906169E42SLOT45END.

Combining WS6--WS9 proves the complete spectra, including multiplicities:

ENDPOINTSOURCEB2396A45C7A5608906169E42SLOT46END

Both traces are ENDPOINTSOURCEB2396A45C7A5608906169E42SLOT47END. Both eigenvalue-two spaces have dimension zero when ENDPOINTSOURCEB2396A45C7A5608906169E42SLOT48END; no nonexistent vector or inverse is used.

Equivalently, introduce the actual spectral projections

ENDPOINTSOURCEB2396A45C7A5608906169E42SLOT49END

They satisfy ENDPOINTSOURCEB2396A45C7A5608906169E42SLOT50END, with ranks ENDPOINTSOURCEB2396A45C7A5608906169E42SLOT51END and ENDPOINTSOURCEB2396A45C7A5608906169E42SLOT52END, respectively, and

ENDPOINTSOURCEB2396A45C7A5608906169E42SLOT53END

This is also a useful direct explanation of the repeated eigenvalues in WS10.

\subsubsection{3. Direct proof of the spectral trace extremum}\label{endpoint-window_review--direct-proof-of-the-spectral-trace-extremum}

Let ENDPOINTSOURCEB2396A45C7A5608906169E42SLOT54END be any self-adjoint endomorphism of the same finite-dimensional Hilbert space, with eigenvalues

ENDPOINTSOURCEB2396A45C7A5608906169E42SLOT55END

All repetitions are retained. By the finite spectral theorem there is an orthonormal eigenbasis ENDPOINTSOURCEB2396A45C7A5608906169E42SLOT56END for this \textbf{same} source form. This use of coordinates is an isometry of the original form, not a change of any vector's norm or any scalar mass.

For an orthogonal projection ENDPOINTSOURCEB2396A45C7A5608906169E42SLOT57END of rank ENDPOINTSOURCEB2396A45C7A5608906169E42SLOT58END, put ENDPOINTSOURCEB2396A45C7A5608906169E42SLOT59END. Then

ENDPOINTSOURCEB2396A45C7A5608906169E42SLOT60END

For ENDPOINTSOURCEB2396A45C7A5608906169E42SLOT61END, set ENDPOINTSOURCEB2396A45C7A5608906169E42SLOT62END. The difference between the largest-ENDPOINTSOURCEB2396A45C7A5608906169E42SLOT63END eigenvalue sum and the last expression is exactly

ENDPOINTSOURCEB2396A45C7A5608906169E42SLOT64END

To verify equality, expand the right side; its coefficient of ENDPOINTSOURCEB2396A45C7A5608906169E42SLOT65END is ENDPOINTSOURCEB2396A45C7A5608906169E42SLOT66END, leaving the displayed difference. Each remaining summand is nonnegative because of the sorted eigenvalues and WS14. Applying this calculation to ENDPOINTSOURCEB2396A45C7A5608906169E42SLOT67END gives the lower bound. When ENDPOINTSOURCEB2396A45C7A5608906169E42SLOT68END the trace and both sums are zero; when ENDPOINTSOURCEB2396A45C7A5608906169E42SLOT69END, ENDPOINTSOURCEB2396A45C7A5608906169E42SLOT70END and both sides equal ENDPOINTSOURCEB2396A45C7A5608906169E42SLOT71END. Thus in every case

ENDPOINTSOURCEB2396A45C7A5608906169E42SLOT72END

This proof requires no strict eigenvalue inequality. It also gives the precise equality condition. At the upper endpoint, ENDPOINTSOURCEB2396A45C7A5608906169E42SLOT73END is identity on all eigenspaces with eigenvalues strictly above the cutoff ENDPOINTSOURCEB2396A45C7A5608906169E42SLOT74END, zero on all eigenspaces strictly below it, and may select the needed rank within the cutoff eigenspace. Indeed a strict coefficient in WS15 forces ENDPOINTSOURCEB2396A45C7A5608906169E42SLOT75END or ENDPOINTSOURCEB2396A45C7A5608906169E42SLOT76END, and those identities imply ENDPOINTSOURCEB2396A45C7A5608906169E42SLOT77END or ENDPOINTSOURCEB2396A45C7A5608906169E42SLOT78END by orthogonal projection. The lower endpoint has the reversed condition. Repeated cutoff values therefore allow their full indicated subspace choices rather than requiring a distinguished eigenbasis there.

Apply WS16 to each projection in WS11. Their ranks give, for ENDPOINTSOURCEB2396A45C7A5608906169E42SLOT79END or ENDPOINTSOURCEB2396A45C7A5608906169E42SLOT80END,

ENDPOINTSOURCEB2396A45C7A5608906169E42SLOT81END

where the exact endpoints are

ENDPOINTSOURCEB2396A45C7A5608906169E42SLOT82END

The sums of length ENDPOINTSOURCEB2396A45C7A5608906169E42SLOT83END are empty for ENDPOINTSOURCEB2396A45C7A5608906169E42SLOT84END. The upper bound is the sum of the largest ENDPOINTSOURCEB2396A45C7A5608906169E42SLOT85END eigenvalues and the largest ENDPOINTSOURCEB2396A45C7A5608906169E42SLOT86END eigenvalues; the lower bound is the corresponding sum of the smallest eigenvalues. They are attained by nested spectral subspaces of these ranks, so the rank bounds are simultaneously consistent with the nested form WS12.

Subtracting WS17 for ENDPOINTSOURCEB2396A45C7A5608906169E42SLOT87END and ENDPOINTSOURCEB2396A45C7A5608906169E42SLOT88END, in both orders, proves

ENDPOINTSOURCEB2396A45C7A5608906169E42SLOT89END

No commutation between ENDPOINTSOURCEB2396A45C7A5608906169E42SLOT90END, ENDPOINTSOURCEB2396A45C7A5608906169E42SLOT91END, and ENDPOINTSOURCEB2396A45C7A5608906169E42SLOT92END has been assumed. All traces are real: for self-adjoint ENDPOINTSOURCEB2396A45C7A5608906169E42SLOT93END, conjugating ENDPOINTSOURCEB2396A45C7A5608906169E42SLOT94END gives ENDPOINTSOURCEB2396A45C7A5608906169E42SLOT95END.

The bound in WS19 is the optimal one from the separate spectra WS10 and WS13. To prove this assertion explicitly, take the two endomorphisms whose diagonal entries in the eigenbasis of ENDPOINTSOURCEB2396A45C7A5608906169E42SLOT96END are the list

ENDPOINTSOURCEB2396A45C7A5608906169E42SLOT97END

Their spectra are precisely WS10. The first attains ENDPOINTSOURCEB2396A45C7A5608906169E42SLOT98END, the second ENDPOINTSOURCEB2396A45C7A5608906169E42SLOT99END, and their difference attains the positive right side of WS19. Exchanging them attains the negative endpoint. This remains true for repeated ENDPOINTSOURCEB2396A45C7A5608906169E42SLOT100END, with equality inside any repeated block understood through WS15. It proves a sharp statement about the indicated unitary orbits; additional alignment information from the literal polynomial and relation flags remains available and can improve a particular instance.

\subsubsection{4. Integration along the unchanged arithmetic Gram path}\label{endpoint-window_review--integration-along-the-unchanged-arithmetic-gram-path}

The original ENDPOINTSOURCEB2396A45C7A5608906169E42SLOT101END is positive and self-adjoint for ENDPOINTSOURCEB2396A45C7A5608906169E42SLOT102END, because ENDPOINTSOURCEB2396A45C7A5608906169E42SLOT103END is Hermitian and its quadratic form is positive. Thus it is diagonalizable with positive eigenvalues

ENDPOINTSOURCEB2396A45C7A5608906169E42SLOT104END

Factor WS2 in its actual order:

ENDPOINTSOURCEB2396A45C7A5608906169E42SLOT105END

The denominator is invertible because each of its eigenvalues ENDPOINTSOURCEB2396A45C7A5608906169E42SLOT106END is positive. Also ENDPOINTSOURCEB2396A45C7A5608906169E42SLOT107END is Hermitian, proving self-adjointness for the exact source metric used in WS4. Functional calculus on the diagonalization of ENDPOINTSOURCEB2396A45C7A5608906169E42SLOT108END gives its eigenvalues

ENDPOINTSOURCEB2396A45C7A5608906169E42SLOT109END

The derivative of ENDPOINTSOURCEB2396A45C7A5608906169E42SLOT110END with respect to ENDPOINTSOURCEB2396A45C7A5608906169E42SLOT111END is exactly ENDPOINTSOURCEB2396A45C7A5608906169E42SLOT112END. Hence the same ordered list WS21 supplies the ordered list WS13 at every ENDPOINTSOURCEB2396A45C7A5608906169E42SLOT113END, including ENDPOINTSOURCEB2396A45C7A5608906169E42SLOT114END, repeated eigenvalues, and values ENDPOINTSOURCEB2396A45C7A5608906169E42SLOT115END. No eigenvalue branch relabelling is needed.

The original determinant transport proves, with these exact projections,

ENDPOINTSOURCEB2396A45C7A5608906169E42SLOT116END

Here ENDPOINTSOURCEB2396A45C7A5608906169E42SLOT117END is the determinant of the original least-norm quotient Gram, so WS24 is the signed arithmetic/reference endpoint correction. The identity was read in the complete sealed source; this review supplies the new spectral refinement. Its integrand is continuous: projection coefficients are rational expressions in positive finite Grams and their positive relation restrictions, so their inverses remain continuous on the compact path.

The elementary derivative

ENDPOINTSOURCEB2396A45C7A5608906169E42SLOT118END

holds also when ENDPOINTSOURCEB2396A45C7A5608906169E42SLOT119END, where both sides are zero. Thus ENDPOINTSOURCEB2396A45C7A5608906169E42SLOT120END. Apply WS19 at every ENDPOINTSOURCEB2396A45C7A5608906169E42SLOT121END, integrate, and use the triangle inequality in WS24. This gives the complete ordered-eigenvalue estimate

ENDPOINTSOURCEB2396A45C7A5608906169E42SLOT122END

All ratios in WS25 are at least one and are literal ratios of eigenvalues of the original relative Gram. Every such ratio is at most ENDPOINTSOURCEB2396A45C7A5608906169E42SLOT123END. The total positive coefficient count is ENDPOINTSOURCEB2396A45C7A5608906169E42SLOT124END. Therefore

ENDPOINTSOURCEB2396A45C7A5608906169E42SLOT125END

For ENDPOINTSOURCEB2396A45C7A5608906169E42SLOT126END, the statement is exactly ENDPOINTSOURCEB2396A45C7A5608906169E42SLOT127END; there is no sum and no factor two. If all ENDPOINTSOURCEB2396A45C7A5608906169E42SLOT128END equal a positive scalar ENDPOINTSOURCEB2396A45C7A5608906169E42SLOT129END, then WS23 makes ENDPOINTSOURCEB2396A45C7A5608906169E42SLOT130END a scalar multiple of the identity and the traces of ENDPOINTSOURCEB2396A45C7A5608906169E42SLOT131END cancel. Equations WS25--WS26 give zero as well. The original common scalar mass remains present in the individual Grams; it cancels only in this signed ratio.

The coefficient ENDPOINTSOURCEB2396A45C7A5608906169E42SLOT132END cannot be decreased using only the separate spectra of ENDPOINTSOURCEB2396A45C7A5608906169E42SLOT133END and the extremes of ENDPOINTSOURCEB2396A45C7A5608906169E42SLOT134END. For any retained numbers ENDPOINTSOURCEB2396A45C7A5608906169E42SLOT135END, choose the ordered data

ENDPOINTSOURCEB2396A45C7A5608906169E42SLOT136END

Then every ratio in WS25 is ENDPOINTSOURCEB2396A45C7A5608906169E42SLOT137END, so ENDPOINTSOURCEB2396A45C7A5608906169E42SLOT138END. Choose the two commuting spectral flags WS20 in a fixed ENDPOINTSOURCEB2396A45C7A5608906169E42SLOT139END-orthogonal eigenbasis of ENDPOINTSOURCEB2396A45C7A5608906169E42SLOT140END. Since ENDPOINTSOURCEB2396A45C7A5608906169E42SLOT141END only changes each basis vector's squared norm by the positive factor ENDPOINTSOURCEB2396A45C7A5608906169E42SLOT142END, these coordinate projections stay orthogonal for every original ENDPOINTSOURCEB2396A45C7A5608906169E42SLOT143END in this abstract comparison. They attain WS19 with the same sign at every ENDPOINTSOURCEB2396A45C7A5608906169E42SLOT144END, hence its integrated value is exactly WS25 and WS26. This is an equality construction within the stated spectral-data class, preserving arbitrary positive ENDPOINTSOURCEB2396A45C7A5608906169E42SLOT145END. It does not identify arbitrary spectral flags with the actual multiplication-by-ENDPOINTSOURCEB2396A45C7A5608906169E42SLOT146END flag.

\subsubsection{5. What the literal polynomial flag retains beyond the spectra}\label{endpoint-window_review--what-the-literal-polynomial-flag-retains-beyond-the-spectra}

The original relation and polynomial flags impose the exact incidence

ENDPOINTSOURCEB2396A45C7A5608906169E42SLOT147END

A nonzero product ENDPOINTSOURCEB2396A45C7A5608906169E42SLOT148END has degree ENDPOINTSOURCEB2396A45C7A5608906169E42SLOT149END, whereas every polynomial on the right has degree at most ENDPOINTSOURCEB2396A45C7A5608906169E42SLOT150END; this proves WS28. In particular independent opposite eigenvalue alignments need not be attainable by these specific subspaces.

For example, assume ENDPOINTSOURCEB2396A45C7A5608906169E42SLOT151END has strictly increasing eigenvalues and ENDPOINTSOURCEB2396A45C7A5608906169E42SLOT152END. Equality at the positive endpoint of WS19 forces ENDPOINTSOURCEB2396A45C7A5608906169E42SLOT153END to have its eigenvalue-two space on the largest ENDPOINTSOURCEB2396A45C7A5608906169E42SLOT154END eigenspaces of ENDPOINTSOURCEB2396A45C7A5608906169E42SLOT155END, and forces ENDPOINTSOURCEB2396A45C7A5608906169E42SLOT156END to vanish on its largest ENDPOINTSOURCEB2396A45C7A5608906169E42SLOT157END eigenspaces, by the equality criterion in WS15. Hence ENDPOINTSOURCEB2396A45C7A5608906169E42SLOT158END. But ENDPOINTSOURCEB2396A45C7A5608906169E42SLOT159END and has positive dimension ENDPOINTSOURCEB2396A45C7A5608906169E42SLOT160END, contradicting WS28. At the negative endpoint the same argument uses the smallest eigenspaces: minimizing ENDPOINTSOURCEB2396A45C7A5608906169E42SLOT161END puts its eigenvalue-two space in the smallest ENDPOINTSOURCEB2396A45C7A5608906169E42SLOT162END eigenspaces, while maximizing ENDPOINTSOURCEB2396A45C7A5608906169E42SLOT163END makes those eigenspaces part of its kernel. Thus that endpoint is also impossible for these literal flags and simple ENDPOINTSOURCEB2396A45C7A5608906169E42SLOT164END.

At ENDPOINTSOURCEB2396A45C7A5608906169E42SLOT165END, the actual formulas reduce to

ENDPOINTSOURCEB2396A45C7A5608906169E42SLOT166END

For simple ENDPOINTSOURCEB2396A45C7A5608906169E42SLOT167END, a positive equality in WS19 makes ENDPOINTSOURCEB2396A45C7A5608906169E42SLOT168END the lowest eigenline and ENDPOINTSOURCEB2396A45C7A5608906169E42SLOT169END the highest eigenline. They are orthogonal, which would force ENDPOINTSOURCEB2396A45C7A5608906169E42SLOT170END. This is impossible because the nonzero constant polynomial is not divisible by the degree-one ENDPOINTSOURCEB2396A45C7A5608906169E42SLOT171END. The negative endpoint is excluded by exchanging the lowest and highest eigenlines. The strictness statement concerns nonconstant simple spectral data; repeated cutoff eigenspaces have precisely the freedom retained in WS15.

These incidence calculations identify an exact additional relation that the source-specific unitary, commutator, and full-jet observations can use. The spectral estimate WS25 remains valid for the original flags, while its equality statement is explicitly scoped. The signed determinant correction itself remains WS24, with all four original endpoints and all relation primitives present.

\subsubsection{6. The original one-dimensional quotient gives a nonlinear bound}\label{endpoint-window_review--the-original-one-dimensional-quotient-gives-a-nonlinear-bound}

This section concerns exactly ENDPOINTSOURCEB2396A45C7A5608906169E42SLOT172END, with the original common source ENDPOINTSOURCEB2396A45C7A5608906169E42SLOT173END, original vertical line ENDPOINTSOURCEB2396A45C7A5608906169E42SLOT174END, original positive measures ENDPOINTSOURCEB2396A45C7A5608906169E42SLOT175END, and original degree-one relation polynomial ENDPOINTSOURCEB2396A45C7A5608906169E42SLOT176END. It makes no assertion that the same two-line formula applies in higher quotient dimension.

The four cutoff indices become ENDPOINTSOURCEB2396A45C7A5608906169E42SLOT177END, so the two copies of ENDPOINTSOURCEB2396A45C7A5608906169E42SLOT178END cancel in this specified ratio and

ENDPOINTSOURCEB2396A45C7A5608906169E42SLOT179END

Let ENDPOINTSOURCEB2396A45C7A5608906169E42SLOT180END be the original constant polynomial, ENDPOINTSOURCEB2396A45C7A5608906169E42SLOT181END, and ENDPOINTSOURCEB2396A45C7A5608906169E42SLOT182END its canonical minimum-norm representative. Every polynomial with remainder one differs from ENDPOINTSOURCEB2396A45C7A5608906169E42SLOT183END by an element of ENDPOINTSOURCEB2396A45C7A5608906169E42SLOT184END, so orthogonal minimization gives ENDPOINTSOURCEB2396A45C7A5608906169E42SLOT185END. The norms are the unchanged quotient values

ENDPOINTSOURCEB2396A45C7A5608906169E42SLOT186END

The last equality follows by the orthogonal splitting of ENDPOINTSOURCEB2396A45C7A5608906169E42SLOT187END; its sign and phase are fixed. Positivity of ENDPOINTSOURCEB2396A45C7A5608906169E42SLOT188END follows independently from ENDPOINTSOURCEB2396A45C7A5608906169E42SLOT189END, which makes ENDPOINTSOURCEB2396A45C7A5608906169E42SLOT190END. Hence the angle between the two lines has

ENDPOINTSOURCEB2396A45C7A5608906169E42SLOT191END

This ratio is strictly below one for the actual source measures. Indeed equality would give ENDPOINTSOURCEB2396A45C7A5608906169E42SLOT192END, hence both ENDPOINTSOURCEB2396A45C7A5608906169E42SLOT193END and ENDPOINTSOURCEB2396A45C7A5608906169E42SLOT194END. On the retained vertical line, ENDPOINTSOURCEB2396A45C7A5608906169E42SLOT195END, so

ENDPOINTSOURCEB2396A45C7A5608906169E42SLOT196END

The two vanishing integrals would force ENDPOINTSOURCEB2396A45C7A5608906169E42SLOT197END. This contradicts positive definiteness of the original Gram on the nonzero polynomial ENDPOINTSOURCEB2396A45C7A5608906169E42SLOT198END; alternatively positivity almost everywhere of ENDPOINTSOURCEB2396A45C7A5608906169E42SLOT199END and the finite zero set of ENDPOINTSOURCEB2396A45C7A5608906169E42SLOT200END give the contradiction directly. Thus ENDPOINTSOURCEB2396A45C7A5608906169E42SLOT201END and ENDPOINTSOURCEB2396A45C7A5608906169E42SLOT202END on the entire compact interval. This argument uses the actual source line and its measure, not an arbitrary coefficient Gram calibration.

Put ENDPOINTSOURCEB2396A45C7A5608906169E42SLOT203END and ENDPOINTSOURCEB2396A45C7A5608906169E42SLOT204END. Their inner product is the positive real number ENDPOINTSOURCEB2396A45C7A5608906169E42SLOT205END. Set ENDPOINTSOURCEB2396A45C7A5608906169E42SLOT206END, and ENDPOINTSOURCEB2396A45C7A5608906169E42SLOT207END. Directly ENDPOINTSOURCEB2396A45C7A5608906169E42SLOT208END and ENDPOINTSOURCEB2396A45C7A5608906169E42SLOT209END. These are explicit coordinates with the original norms displayed, so no Gram factor or mass is removed. In the basis ENDPOINTSOURCEB2396A45C7A5608906169E42SLOT210END, the difference in WS29 is

ENDPOINTSOURCEB2396A45C7A5608906169E42SLOT211END

and is zero on its orthogonal complement. The displayed matrix has trace zero and determinant ENDPOINTSOURCEB2396A45C7A5608906169E42SLOT212END, so the full three eigenvalues are ENDPOINTSOURCEB2396A45C7A5608906169E42SLOT213END. This proves the exact alignment-dependent spectrum

ENDPOINTSOURCEB2396A45C7A5608906169E42SLOT214END

Write its positive and negative rank-one spectral projections as ENDPOINTSOURCEB2396A45C7A5608906169E42SLOT215END and ENDPOINTSOURCEB2396A45C7A5608906169E42SLOT216END. Then ENDPOINTSOURCEB2396A45C7A5608906169E42SLOT217END. Applying the already proved rank-one bound WS16 to their two traces gives

ENDPOINTSOURCEB2396A45C7A5608906169E42SLOT218END

Here ENDPOINTSOURCEB2396A45C7A5608906169E42SLOT219END is the original derivative in WS24; the equality identifying it uses all four original window signs even though their scalar endpoint expression has the cancellation in WS30.

For real ENDPOINTSOURCEB2396A45C7A5608906169E42SLOT220END, let ENDPOINTSOURCEB2396A45C7A5608906169E42SLOT221END, using the nonnegative real inverse of ENDPOINTSOURCEB2396A45C7A5608906169E42SLOT222END. Differentiating ENDPOINTSOURCEB2396A45C7A5608906169E42SLOT223END gives

ENDPOINTSOURCEB2396A45C7A5608906169E42SLOT224END

Since ENDPOINTSOURCEB2396A45C7A5608906169E42SLOT225END is continuous and strictly positive on the compact interval, it has a positive minimum; every composition and division in WS37 is therefore defined and smooth on this whole path. Multiplying WS36 by WS37 yields ENDPOINTSOURCEB2396A45C7A5608906169E42SLOT226END. Integrating WS23 gives the exact nonlinear endpoint control

ENDPOINTSOURCEB2396A45C7A5608906169E42SLOT227END

Equivalently, put ENDPOINTSOURCEB2396A45C7A5608906169E42SLOT228END and ENDPOINTSOURCEB2396A45C7A5608906169E42SLOT229END. Then ENDPOINTSOURCEB2396A45C7A5608906169E42SLOT230END lies between ENDPOINTSOURCEB2396A45C7A5608906169E42SLOT231END and ENDPOINTSOURCEB2396A45C7A5608906169E42SLOT232END. The function ENDPOINTSOURCEB2396A45C7A5608906169E42SLOT233END is increasing for ENDPOINTSOURCEB2396A45C7A5608906169E42SLOT234END, because its derivative is ENDPOINTSOURCEB2396A45C7A5608906169E42SLOT235END. Hence

ENDPOINTSOURCEB2396A45C7A5608906169E42SLOT236END

The lower bound may be zero although the actual value was proved strictly positive. A scalar relative Gram gives ENDPOINTSOURCEB2396A45C7A5608906169E42SLOT237END and the exact equality of the two original endpoint values, in agreement with WS25. This is an additional analytic calculation on the actual ENDPOINTSOURCEB2396A45C7A5608906169E42SLOT238END source, retaining its canonical line angle.

\subsubsection{7. Complete independent review of the root AW1--AW23 source}\label{endpoint-window_review--complete-independent-review-of-the-root-aw1aw23-source}

The complete \texttt{tau\_relation\_window\_spectral\_transport\_20260913.tex} was read and checked at SHA-256 \texttt{c79e76c087efd59871e3234ff44fb18964515dbd34f0318dffba6b3f1440f6ec}. No mathematical correction is requested at that pin.

The following checks cover the new assertions rather than only their final inequalities.

\begin{enumerate}
\def\labelenumi{\arabic{enumi}.}
\tightlist
\item
  \textbf{Original objects and types.} AW1--AW6 retain the exact arithmetic Taylor unit, two distinct injected observations into the full tensor jet module and into ENDPOINTSOURCEB2396A45C7A5608906169E42SLOT239END, the common ENDPOINTSOURCEB2396A45C7A5608906169E42SLOT240END coefficient cutoff, the full masses, and the same-metric projection formulas. The correction sign agrees with the four terms of the sealed AT15, whose source pin is \texttt{24cf7ce5be6a43cc656641b516185df10b2e4d7824600ece657120d1c549deae}.
\item
  \textbf{Smooth actual isometry.} The Gram recursion AW9 retains every positive norm ENDPOINTSOURCEB2396A45C7A5608906169E42SLOT241END. Linear independence of each original list makes every residual nonzero; positive square roots are smooth along the entire compact path. The canonical quotient list lies in ENDPOINTSOURCEB2396A45C7A5608906169E42SLOT242END because of the original minimum-section orthogonality and spans it because ENDPOINTSOURCEB2396A45C7A5608906169E42SLOT243END. Thus AW10 maps an actual orthonormal basis to another one, with matching eigenvalues from WS7 and WS9. It proves ENDPOINTSOURCEB2396A45C7A5608906169E42SLOT244END and ENDPOINTSOURCEB2396A45C7A5608906169E42SLOT245END, including the empty middle list at ENDPOINTSOURCEB2396A45C7A5608906169E42SLOT246END.
\item
  \textbf{Commutator and unchanged cohomology observation.} Finite trace cyclicity gives ENDPOINTSOURCEB2396A45C7A5608906169E42SLOT247END, with precisely the displayed sign. The original observation has kernel ENDPOINTSOURCEB2396A45C7A5608906169E42SLOT248END because its two later maps are injective. Therefore ENDPOINTSOURCEB2396A45C7A5608906169E42SLOT249END. The explicit AW10 basis assignment identifies this with ENDPOINTSOURCEB2396A45C7A5608906169E42SLOT250END. The full observation defect ENDPOINTSOURCEB2396A45C7A5608906169E42SLOT251END is retained, so this isometry is not assigned an unwritten identity on arithmetic classes. The support lift preserves the label and gives exactly that supported-zero inverse image.
\item
  \textbf{Repeated spectra and integration.} AW13--AW17 agree with the independently derived WS14--WS26. The zero rank at ENDPOINTSOURCEB2396A45C7A5608906169E42SLOT252END, repeated Gram eigenvalues, values ENDPOINTSOURCEB2396A45C7A5608906169E42SLOT253END, and the full coefficient ENDPOINTSOURCEB2396A45C7A5608906169E42SLOT254END are all covered. The claimed sharpness is expressly restricted to the separate spectral constraint; its text does not freely orient the original polynomial relations. AW18 is a correct substitution of the supplied Gamma lower estimate into this new finite-source bound, with the actual quartet restriction retained.
\item
  \textbf{One-dimensional continuation.} AW19--AW23 are independently proved in full by WS30--WS39 above. The vertical-line identity, positive phase of ENDPOINTSOURCEB2396A45C7A5608906169E42SLOT255END, exact two-dimensional projection matrix, remaining zero eigenvalue, positive derivative divisor, and both real hyperbolic-function branches are correct. The factor on the right of AW22 is ENDPOINTSOURCEB2396A45C7A5608906169E42SLOT256END, while AW23 correctly uses half that number in its ENDPOINTSOURCEB2396A45C7A5608906169E42SLOT257END arguments. The calculation remains explicitly restricted to ENDPOINTSOURCEB2396A45C7A5608906169E42SLOT258END.
\end{enumerate}

The whole AW source is accepted at the exact pin above. The original finite Gram data remain the input to the proved estimates; no uniform bound in ENDPOINTSOURCEB2396A45C7A5608906169E42SLOT259END has been presumed.

\subsubsection{8. Source and verification scope}\label{endpoint-window_review--source-and-verification-scope}

The complete sealed \texttt{tau\_arithmetic\_determinant\_transport\_20260913.tex} was read. Its AT11--AT15 fix the original common Hilbert space, projection conventions, correction sign, and relative Gram path used here. This note proves WS5--WS39 directly from those data and the actual ENDPOINTSOURCEB2396A45C7A5608906169E42SLOT260END measures; it does not alter the sealed source or its PDF. The complete root AW1--AW23 source was also read and accepted as recorded above. These complete written proofs constitute the closed mathematical review; no supplementary source is needed for any acceptance recorded here.

No external trace extremum theorem was used: WS14--WS16 are its direct finite-dimensional proof, including repeated eigenvalues. No Lean execution, remote mutation, numerical-only proof, source-multiplier boundedness hypothesis, or unproved asymptotic estimate is asserted.
